% ============================================================================
% PREAMBLE.tex (REVISED & OPTIMIZED)
% Buku Ajar Pemrograman Game
% ============================================================================

% ============================================================================
% BASIC SETTINGS
% ============================================================================

% Encoding & language
\usepackage[T1]{fontenc}
\usepackage[indonesian]{babel}
% inputenc optional (pdfLaTeX legacy)
\usepackage[utf8]{inputenc} % boleh dihapus jika pakai XeLaTeX/LuaLaTeX

% Font: Times New Roman sesuai standar buku ajar Indonesia
\usepackage{times} % Times New Roman font family

% Geometry - Sesuai standar buku ajar Indonesia
% Margin: atas 2.5cm, kiri 3cm, bawah 2cm, kanan 2cm
\usepackage[a4paper, left=3cm, right=2cm, top=2.5cm, bottom=2cm]{geometry}

% Typography improvement (expansion=false: avoids pdfTeX error with non-scalable fonts)
\usepackage[expansion=false]{microtype}

% ============================================================================
% TABLE PACKAGES
% ============================================================================

\usepackage{booktabs}
\usepackage{longtable}
\usepackage{array}
\usepackage{multirow}
\usepackage{tabularx}
\usepackage{colortbl}

% ============================================================================
% MATH & ALGORITHM
% ============================================================================

\usepackage{amsmath,amssymb,amsthm,mathtools}
\usepackage{algorithm}
\usepackage{algpseudocode}

% === CODE LISTING ===
% Pilih salah satu: listings ATAU minted
\usepackage{listings}
% \usepackage{minted} % aktifkan jika pakai -shell-escape

% ============================================================================
% GRAPHICS & DIAGRAM
% ============================================================================

\usepackage{graphicx}
\usepackage{tikz}
\usepackage{pgfplots}
\pgfplotsset{compat=1.18}
\usetikzlibrary{shapes,arrows,positioning,calc,decorations.pathmorphing}
% TikZ styles for block diagrams (used in ch01, ch07, etc.)
\tikzset{
  block/.style={rectangle, draw, fill=blue!20, minimum height=1.2em, minimum width=2.5em, align=center, rounded corners=2pt},
  line/.style={draw}
}
\usepackage{float}

% ============================================================================
% REFERENCES & BIBLIOGRAPHY
% ============================================================================

\usepackage[autostyle=false,style=fallback]{csquotes}
\usepackage[style=apa,backend=biber,language=english]{biblatex}
\DeclareLanguageMapping{indonesian}{english-apa}
\DeclareLanguageMapping{english}{english-apa}
\addbibresource{references.bib}

% ============================================================================
% LISTS & STRUCTURE
% ============================================================================

\usepackage{enumitem}
\setlist[itemize]{leftmargin=*,topsep=0.5em}
\setlist[enumerate]{leftmargin=*,topsep=0.5em}

% === CUSTOM OBE ENVIRONMENTS (Learning Outcomes, Definisi, Contoh, Latihan, Catatan) ===
\newtheorem{definisi}{Definisi}[section]
\newtheorem{contoh}{Contoh}[section]
\newtheorem{latihan}{Latihan}[section]
\newtheorem{catatan}{Catatan}[section]
\newtheorem{tips}{Tips}[section]
\newenvironment{learningoutcomes}{%
  \par\noindent\textbf{Learning Outcomes:}\par\vspace{0.3em}%
  \begin{itemize}[leftmargin=*, itemsep=0.2em]%
}{%
  \end{itemize}%
}

\usepackage{subfiles}

% ============================================================================
% UTILITIES
% ============================================================================

\usepackage{xcolor}
\usepackage{fancyhdr}
\usepackage{titlesec}
\usepackage{tocloft}
\PassOptionsToPackage{nowarn}{tracklang}
\usepackage[nogroupskip,nolangwarn]{glossaries}
\usepackage{makeidx}
\usepackage{url}
\usepackage{verbatim}

% ============================================================================
% HYPERREF (load near the end)
% ============================================================================

\usepackage{hyperref}
\hypersetup{
    colorlinks=true,
    linkcolor=blue,
    urlcolor=cyan,
    citecolor=red,
    pdftitle={Buku Ajar Pemrograman Game},
    pdfauthor={Program Studi S1 Informatika},
    pdfsubject={Game Development dengan Unity},
    pdfkeywords={Unity, C\#, Game Development, OBE}
}

% Optional: smarter cross-references
\usepackage[nameinlink,noabbrev]{cleveref}

% ============================================================================
% NUMBERING CONFIGURATION
% ============================================================================

\renewcommand{\thechapter}{\arabic{chapter}}
\renewcommand{\thesection}{\thechapter.\arabic{section}}
\renewcommand{\thesubsection}{\thesection.\arabic{subsection}}

\counterwithin{figure}{chapter}
\counterwithin{table}{chapter}
\counterwithin{equation}{chapter}

% ============================================================================
% LISTINGS CONFIG (C#)
% ============================================================================

\lstset{
    language=[Sharp]C,
    basicstyle=\ttfamily\small,
    keywordstyle=\color{blue},
    commentstyle=\color{green!60!black},
    stringstyle=\color{red},
    numbers=left,
    numberstyle=\tiny\color{gray},
    backgroundcolor=\color{gray!10},
    frame=single,
    breaklines=true,
    tabsize=2
}

% ============================================================================
% HEADER & FOOTER
% ============================================================================

\setlength{\headheight}{14.5pt}
\pagestyle{fancy}
\fancyhf{}
\fancyhead[LE]{\leftmark}
\fancyhead[RO]{\rightmark}
\fancyfoot[C]{\thepage}

% ============================================================================
% CHAPTER TITLE FORMAT
% Sesuai standar buku ajar: Judul bab 16pt, subjudul 14pt
% ============================================================================

\titleformat{\chapter}[display]
{\normalfont\fontsize{16}{20}\bfseries\selectfont}
{\chaptertitlename\ \thechapter}{20pt}{\fontsize{16}{20}\bfseries\selectfont}
\titlespacing*{\chapter}{0pt}{50pt}{40pt}

% Format section (subjudul) 14pt
\titleformat{\section}
{\normalfont\fontsize{14}{18}\bfseries\selectfont}
{\thesection}{1em}{}

% Format subsection 12pt bold
\titleformat{\subsection}
{\normalfont\fontsize{12}{15}\bfseries\selectfont}
{\thesubsection}{1em}{}

% ============================================================================
% GLOSSARY & INDEX
% ============================================================================

\makeglossaries
\setacronymstyle{long-short}
\loadglsentries{glossary}

\makeindex

% ============================================================================
% SPACING
% Sesuai standar buku ajar: spasi 1.5, font 12pt
% ============================================================================

\setlength{\parindent}{1.5em}
\setlength{\parskip}{0.5em}
\renewcommand{\baselinestretch}{1.5} % Spasi 1.5 sesuai standar

\raggedbottom
\sloppy
\emergencystretch=3em
